\section{A connectivity test from $A$: the answer to question 2}\label{sec:connectivity}

A positive $(A^k)_{ij}$ says by \Cref{thm:walks} only that a walk of length $k$ exists, and a walk may revisit vertices, while \Cref{def:connected} asks for paths. The next lemma closes that gap and also bounds the lengths to be checked.

\plainstart\textbf{In plain words.} Any stroll between two houses can be cut down to a no-return stroll between them. Whenever the stroll visits a house twice, cut out the part in between. A no-return stroll takes fewer steps than the village has houses.\plainend
\begin{lemma}[Walk contains path]\label{lem:walkpath}
Let $G$ have $N$ vertices.
(a) If there is a walk of length $k$ from $n_i$ to $n_j$ ($i = j$ allowed), there is a path from $n_i$ to $n_j$ of length at most $k$ whose vertices all occur in the walk.
(b) A path in $G$ has length at most $N - 1$.
\end{lemma}

Grinberg proves (a) without the clause that the path's vertices occur in the walk and states (b) \cite[Corollary 2.9.7, p.~39; p.~40]{grinberg_graph_theory}; the textbook's Lemma 1.2.1 is stated only for two different vertices and has no length bound \cite[p.~23]{petrosyan_graph_theory}. The proofs below are this paper's own.

\begin{proof}
(a) By strong induction on $k$. \emph{Base case} $k = 0$: the walk is $(n_i)$, so $n_j = n_i$, and $(n_i)$ has no repeated vertex, so it is a path of length $0$. \emph{Step:} let $k \ge 1$ and assume (a) for every walk of smaller length. Let $w = (w_0, \dots, w_k)$ be a walk from $n_i$ to $n_j$. If no vertex repeats in $w$, then $w$ is a path of length $k$. Otherwise $w_a = w_b$ for some $a < b$. Form $w' = (w_0, \dots, w_a, w_{b+1}, \dots, w_k)$, which is $(w_0, \dots, w_a)$ when $b = k$. It is a walk: consecutive terms inside the two kept blocks were adjacent in $w$, and the one new consecutive pair is $(w_a, w_{b+1})$, which equals $(w_b, w_{b+1})$ because $w_a = w_b$, and so is adjacent in $w$. It starts at $w_0 = n_i$, and it ends at $w_k = n_j$ (when $b = k$ it ends at $w_a = w_b = w_k$). Its length is $k - (b - a) < k$. By the induction hypothesis $w'$ contains a path from $n_i$ to $n_j$ of length at most $k - (b-a) \le k$, with vertices occurring in $w'$ and so in $w$.

(b) A path of length $k$ has $k + 1$ terms $w_0, \dots, w_k$, pairwise distinct vertices of $G$. $G$ has $N$ vertices, so $k + 1 \le N$.
\end{proof}

\plainstart\textbf{In plain words.} The textbook calls a village joined-up using a weaker kind of stroll, one that never reuses a lane. A no-return stroll never reuses a lane. Any stroll that never reuses a lane can be cut down to a no-return stroll, as shown below.\plainend
\begin{proposition}[Connectedness agrees with the textbook's definition]\label{prop:same}
For all $n_i, n_j$ in a simple graph, there is a path from $n_i$ to $n_j$ (\Cref{def:path}) if and only if there is a trail from $n_i$ to $n_j$. Hence the graphs that are connected by \Cref{def:connected} are the graphs that are connected by the textbook's Definition 2.1.1, which uses its own ``path'', the trail of \Cref{rem:trap}.
\end{proposition}

\begin{proof}
\emph{Path implies trail.} Let $(v_0, \dots, v_k)$ be a path and suppose two of its edges are equal, $\{v_{s-1}, v_s\} = \{v_{t-1}, v_t\}$ with $s < t$. Then $v_{s-1} = v_{t-1}$ or $v_{s-1} = v_t$. In the first case positions $s-1$ and $t-1$ hold the same vertex, in the second positions $s-1$ and $t$ do; these positions differ because $s - 1 < t - 1 < t$. Either way a vertex of the path repeats, a contradiction. So the edges are pairwise distinct and the path is a trail. \emph{Trail implies path.} A trail from $n_i$ to $n_j$ is a walk from $n_i$ to $n_j$, so by \Cref{lem:walkpath}(a) it contains a path from $n_i$ to $n_j$. The two directions give the equivalence for each pair, hence for the quantified definitions.
\end{proof}

\plainstart\textbf{In plain words.} Chain the neighbour chart for one step, two steps, and so on. Stop at one less than the number of houses, and add the charts box by box. The village is joined-up exactly when each box for two different houses holds more than zero.\plainend
\begin{theorem}[Sum test]\label{thm:test}
Let $N \ge 2$ and $S = A + A^2 + \dots + A^{N-1}$. The graph is connected if and only if $S_{ij} > 0$ for all $i \ne j$.
\end{theorem}

Paramadevan and Sotheeswaran state this test \cite[Theorem 2, p.~17]{paramadevan2021adjacency}; the proof below is this paper's own.

\begin{proof}
\emph{Step 0.} By \Cref{thm:walks} each $(A^k)_{ij}$ is a number of walks, a non-negative integer. So $S_{ij} = W_1(i,j) + \dots + W_{N-1}(i,j)$, and $S_{ij} > 0$ if and only if for some $k$ with $1 \le k \le N-1$ there is a walk of length $k$ from $n_i$ to $n_j$.

\emph{Connected implies $S_{ij} > 0$.} Let $i \ne j$. By \Cref{def:connected} there is a path $P$ from $n_i$ to $n_j$, of length $k$ say. Since $n_i \ne n_j$, $P$ has at least two vertices, so $k \ge 1$, and $k \le N - 1$ by \Cref{lem:walkpath}(b). $P$ is a walk of length $k$, so $S_{ij} \ge W_k(i,j) \ge 1$.

\emph{$S_{ij} > 0$ for all $i \ne j$ implies connected.} Take any pair $n_i, n_j$. If $i = j$, the path $(n_i)$ of length $0$ serves. If $i \ne j$, Step 0 gives a walk from $n_i$ to $n_j$, and \Cref{lem:walkpath}(a) gives a path from $n_i$ to $n_j$. Every pair has a path, so $G$ is connected.
\end{proof}

The bound $N-1$ is used only in the first direction. The diagonal is left out because it is not needed (each vertex has its length-$0$ path) and because it can be zero in a connected graph: for $N = 2$ and one edge, $S = A$ has zero diagonal. For $N = 1$ the graph is connected and no pair has $i \ne j$, so the condition holds vacuously, while $S$ is an empty sum; the test is therefore stated for $N \ge 2$.

\begin{remark}[The form $(I+A)^{N-1}$]\label{rem:binom}
Gotti and Bai state the test with $R = (I + A)^{N-1}$ in place of $S$, where $I$ is the identity matrix \cite[Proposition 3, p.~2]{gotti_18211_lecture27} \cite[Theorem 2, p.~1]{bai_bimsa_adjacency}. Since $I$ and $A$ commute, the binomial theorem gives $R = \sum_{k=0}^{N-1} \binom{N-1}{k} A^k$. Every coefficient is at least $1$ and every $A^k$ has non-negative entries, so for $i \ne j$ the entry $R_{ij}$ is positive exactly when $S_{ij}$ is (the term $k = 0$ adds only to the diagonal), and each $R_{ii} \ge 1$. So $R$ has no zero entry exactly when the test of \Cref{thm:test} holds. By script, $R$ has rows $(2,2,1), (2,3,2), (1,2,2)$ for $G_1$ of \Cref{ex:board} and $(2,2,0), (2,2,0), (0,0,1)$ for $G_2$ of \Cref{ex:test2}.
\end{remark}
